% Authoritative editable manuscript. Generated reading text and PDF use scripts/build_paper.py.
\documentclass[11pt,a4paper]{article}
\usepackage[T1]{fontenc}
\usepackage{lmodern}
\usepackage{amsmath,amssymb}
\usepackage[margin=27mm]{geometry}
\usepackage{microtype}
\usepackage{booktabs,longtable,array}
\usepackage{xurl}
\usepackage{listings}
\lstset{literate={→}{{$\to$}}1 {∀}{{$\forall$}}1 {∃}{{$\exists$}}1 {∧}{{$\land$}}1 {¬}{{$\neg$}}1 {≤}{{$\le$}}1,basicstyle=\ttfamily\footnotesize,columns=fullflexible,keepspaces=true,breaklines=true,frame=single,framerule=0.2pt}
\usepackage[hidelinks,pdfauthor={Roger Gujord},pdftitle={The base-seven Heraclitus transform: computer-assisted coverage}]{hyperref}
\usepackage{fancyhdr}
\pagestyle{fancy}
\fancyhf{}
\fancyhead[L]{\small The base-seven Heraclitus transform}
\fancyhead[R]{\small Version 1.1.0-rc.1}
\fancyfoot[C]{\thepage}
\setlength{\headheight}{14pt}
\setlength{\parskip}{4pt}
\setlength{\parindent}{0pt}
\widowpenalty=10000
\clubpenalty=10000
\displaywidowpenalty=10000
\emergencystretch=1em
\title{The base-seven Heraclitus transform\\[5pt]\large Computer-assisted coverage}
\author{Roger Gujord\\\small ORCID: \href{https://orcid.org/0009-0006-5227-980X}{0009-0006-5227-980X}}
\date{9 November 2026\\\small Local revision 1.1.0-rc.1}
\begin{document}
\maketitle
\begin{abstract}
Starting at zero, the base-seven Heraclitus transform repeatedly chooses an
unused integer at a power-of-seven distance from its current position,
minimizing absolute value and preferring the positive integer in a tie.
We present a proposed computer-assisted proof that this walk visits every
integer exactly once. The argument reduces the infinite walk to a closed
catalogue of normalized boundary states with period 36 across scales.
An exact finite base and a uniform scale transition yield coverage of
$[-C_k+37,C_k-37]$ by index $7^k$, where $C_k=(7^k-1)/2$ and $k\ge6$.
The accompanying certificates have a separate Python checker; a Lean
companion formalizes the exact rule, base, transitions and coverage.
Conditional on the uniform transition, a first-crossing argument gives
$\limsup_{n\to\infty}\bigl||a(n)|-n/2\bigr|/n\ge3/8$.
A Lean extension now establishes the exact rational-epsilon contract for
equality at $3/8$ from the original rule, using the formalized base and
scale transitions. Its upper estimate covers every actual index, with
complete used-set, tail, time and radius bounds throughout each scale.
This AI-assisted manuscript awaits independent human mathematical review.
\end{abstract}

\section*{Introduction}
Greedy walks with restricted jump lengths combine a local choice rule with
a global coverage problem: can a rule that always chooses the unused
integer closest to zero eventually visit every integer? Here the allowed
jump lengths are $1,7,49,\ldots$. The challenge is to control the complete
visited set at arbitrarily large scales, including isolated points marked
by earlier long jumps.

Coverage for powers of two already has a public block description \cite{ref1}.
For powers of seven, the reachable boundary states involve 74- and 130-step
local cycles and a 36-phase description across scales. This makes the base
a concrete setting in which to study a finite-state approach to coverage.
The finite closure established here is specific to seven; no result for
another odd base is assumed.

The central mechanism is a uniform marked-endpoint transfer. Translation
cycles describe long stretches of the walk; a finite repair handles the
next previsited endpoint. Six such transfers form each ladder layer.
Closing both the entire endpoint state and its spacing residue gives an
induction over arbitrarily many layers, and then over successive scales.
A repeating pattern in sampled prefixes alone would not supply this step.

Sections 1--7 give the coverage argument; Section 8 gives the deviation
bounds. The remaining sections describe verification, reproduction,
limitations and related work.


\section{Statement and precise rule}
\label{sec:1}

Set $a(0)=0$. Having chosen distinct integers $a(0),\ldots,a(n)$, choose $a(n+1)$ among the unused integers $v$ for which $|v-a(n)|$ is a nonnegative integral power of 7. Minimize $|v|$, giving positive $v$ priority in a tie.

There is always an available move: a sufficiently large power reaches outside the finite used set. Thus the definition gives an infinite sequence without repetitions.

\noindent\textbf{Theorem 1 (coverage).} Every integer occurs in $a$. More precisely, put $C_k=(7^k-1)/2$. For every $k\ge6$,

\begin{equation*}
[-C_k+37,C_k-37]\subseteq\{a(0),\ldots,a(7^k)\}. \tag{1}\label{eq:coverage}
\end{equation*}

The complete state at $n=7^k$, translated by $C_k$ separately on each side of zero, has period 36 in $k$ for $k\ge6$. The direct initial check also gives the same phase table at $k=4,5$.

For a fixed integer $z$ choose $k$ large enough that $|z|\le C_k-37$. Then (1) gives an occurrence of $z$. Uniqueness follows from the rule. This last implication is elementary. The substantive work is proving the finite state transition uniformly for arbitrarily large $k$.

\section{States and certificates}
\label{sec:2}

A normalized half-line state is $(c,B)$. All coordinates below 0 are used or forbidden, 0 is unused, $B$ is a finite set of positive used coordinates, and the current coordinate is $c$. Thus $c<0$ or $c$ belongs to $B$. After a greedy choice $v$, adjoin $v$, let $d$ be the least missing nonnegative coordinate, and subtract $d$ from all coordinates. The new state is

\begin{equation*}
\bigl(v-d,\{x-d:x\in B\cup\{v\},\ x\ge d\}\bigr). \tag{2}\label{eq:normalize}
\end{equation*}

and the front advances by $d$. Greedy here means the \emph{least} unused half-line neighbor, not least absolute value on the signed line.

The checker verifies a transition by scanning candidate integers $0,1,\ldots,v$ and testing whether the absolute distance is a power of 7 by repeated integer division. Consequently every smaller possible candidate is excluded, including candidates using larger powers. It does not import the generator's neighbor-selection function.

In a signed state at scale $m=7^k$, write $C=(m-1)/2$. A point on side $s\in\{+1,-1\}$ has the form $s(C+x)$. A side record $(f,B)$ means that every positive magnitude below $C+f$ is used, $C+f$ is unused, and the additional used magnitudes are $C+x$ for $x$ in $B$. $B$ contains only values greater than $f$. The current record is $(s,x)$. Zero is always used. These records specify the entire visited set, not just a bounded sample of it.

The accompanying JSON files contain every finite transition needed below, including complete transient and cycle words. They are part of the computer-assisted proof. A separate, self-contained verifier checks the stated identities and greedy choices.

\section{Translation cycles and a uniform marked-point lemma}
\label{sec:3}

\noindent\textbf{Lemma 2 (translation cycle).} Suppose an exact normalized
half-line greedy trajectory returns to its entire state $(c,B)$ after
$P>0$ choices, with total front shift $P$. Let $D_i$ be the front shift
before choice $i$, $v_i$ the chosen coordinate in the then-current
normalization, and $W_i=D_i+v_i$, for $0\le i<P$. Suppose the $W_i$ give
distinct residues modulo $P$. Then successive cycles choose $W_i+qP$,
$q\ge0$, relative to the front at cycle entry. The certificates check all
these finite hypotheses for the cycles used below.

\emph{Proof.} After the first block the prefix below the new front is filled, the finite pattern $B$ is the same, and the current point has the same relative coordinate. Distances and order of candidates are translation invariant. Backward jumps below the front land in the filled prefix. Induction therefore repeats the identical normalized block. This is an infinite statement derived from the finite transition checks.

\noindent\textbf{Lemma 3 (uniform marked-endpoint transfer).} Work relative
to a completed marked endpoint at coordinate 0. Let $(f,c,B)$, with $f>0$,
describe a valid half-line state: every coordinate below $f$ is used,
$f$ is unused, the remaining used points are $f+B$, and the current point
$f+c$ is used. Add a next previsited endpoint $h\ge2401$ and any finite set
of later previsited endpoints, all at least $2h$. Assume that the walk
without these additional endpoints has a checked finite transient followed
by a cycle satisfying Lemma 2. Assume also that the marker insertion and
finite repair described below are checked exact half-line transitions and
satisfy the five bounds below. Then the returned finite endpoint record,
measured relative to $h$, depends only on the initial state and $h$ modulo
$P$. The complete used set consists of that record and its filled prefix,
together with the untouched later endpoints. Those later endpoints do not
affect the transfer choices or its finite output record.

\emph{Proof and finite hypotheses.} Temporarily ignore the additional
endpoints and follow the unmarked half-line walk. Let $f_*$ be the front at
entry to its translation cycle of length $P$.

There is exactly one cycle phase $i$ with

\begin{equation*}
h-f_*-W_i\equiv0\pmod P.
\end{equation*}

If $q=(h-f_*-W_i)/P\ge0$, the unmarked process first selects $h$ at that phase of cycle $q$. Declaring $h$ used from the beginning changes none of the preceding choices: until a removed candidate would have won, removing it does not change the minimum. At this instant insert the marker, without moving the current point or counting a new choice, normalize by (2), and continue until the front is strictly beyond $h$. This defines the next endpoint state.

At a fixed phase $i$, the unmarked normalized state is the same in every
cycle, and the inserted marker has normalized coordinate $v_i$. Thus its
normalized insertion state and finite repair word are independent of $q$;
the repair can be checked once for each input state and spacing residue.

The finite endpoint record depends only on the input state and $h$ modulo
$P$, provided the transient and repair cannot encounter a different future
endpoint. The following bounds are checked for each certificate. They
refer to the auxiliary walk with future markers omitted, inserting $h$
only at the specified phase:

\begin{itemize}
\item all initially used points and transient choices are below 2000 relative to the preceding endpoint;
\item $f_*+W_i<2000$, guaranteeing $q\ge0$ for $h\ge2401$;
\item $\max(W)-\min(W)<2000$;
\item immediately after inserting $h$, all used points above $h$ are less than $h$+2000;
\item every repair choice, and every used point in the resulting endpoint state, is less than $h$+2000.
\end{itemize}

For $h\ge2401$ the initial region is before $h$. Before the marked phase, any cycle choice is at most $h+\max(W)-\min(W)<2h$. The repair also stays below $2h$. Thus no later premarked endpoint can affect any of these choices. This proves the transfer rule for \emph{every} $h\ge2401$ in its residue class, not just the large $h$ used during discovery. Choices below $h$ are unrestricted by these upper bounds and are checked by the exact half-line rule.

There are 119 distinct transfer certificates in the proof. Their cycle lengths belong to $\{1,4,8,74,130\}$. The generator obtains a largest individual threshold of 376; the separate certificate audit uses the more conservative common bounds 2000 and 2401 above.

The last checkpoint need not be the \emph{first} instant all numbers through $h$ are used. It is the instant specified by the marked-phase construction followed by the finite repair. This distinction is necessary because a marker can allow the actual filled front to pass $h$ before the unmarked walk would have selected $h$.

\section{Arbitrarily many levels of a descending ladder}
\label{sec:4}

The local configurations needed for the signed proof have all coordinates below -100 filled and a finite used set in $[-100,100)$. An offset $A$ with $-100\le A<100$ is already used. Start at $A+7^K$, where $K\ge4$, with no other used points above the finite old region.

Above level 4 the greedy descent is forced. From $A+7^{j+1}$ it visits

\begin{equation*}
A+6\cdot7^j,\ A+5\cdot7^j,\ \ldots,\ A+7^j.
\end{equation*}

A jump $7^{j+1}$ would return to the already-used $A$; at intermediate points it goes below the filled prefix. Smaller backward jumps leave a larger candidate than subtraction of $7^j$. The indicated points are unused. Forward jumps are larger. Since $|A|<100$ and $7^j\ge2401$, no ordering ambiguity arises.

Consequently the high descent premarks the ladder endpoints. At the bottom, start at $A+2401$ with the old finite used set. For each of 18 explicit old configurations the checker follows the full greedy word until every coordinate through $A+2401$ is used. It checks that the largest choice is less than $A+2401+2000$. The next higher premarked endpoint is $A+2\cdot2401$, so it could not have affected this finite bottom computation.

Once the endpoint $A+h$ has been completed, with $h=7^j$, six marked points remain in the next layer:

\begin{equation*}
A+2h,\ A+3h,\ \ldots,\ A+7h.
\end{equation*}

Apply the marked-point lemma six times with spacing $h$. This completes the next layer and returns another finite endpoint state. Increase $h$ to $7h$ and repeat. The certificate records this map.

Let $L=\operatorname{lcm}(1,4,8,74,130)=19240$. Every relevant transfer depends only on its state and $h$ modulo a divisor of $L$, and

\begin{equation*}
7^{36}\equiv1\pmod L. \tag{3}\label{eq:period}
\end{equation*}

For 16 of the 18 entrance configurations, the layer states and the needed residues repeat after 2 layers, following at most one transient layer. For the remaining two they repeat after 36 layers, without a transient. Each repeated pair consists of the \emph{entire} endpoint state and $h$ modulo the modulus used by that ladder, either 8 or 19240. The checker verifies all six transfers at each layer and the exact closing equality. The spacing only increases, so the uniform lower bound $h\ge2401$ persists.

\noindent\textbf{Proposition 4 (ladder closure).} For each of the 18
specified entrance configurations, the checked bottom word, the six
transfers per layer and closure of the pair (entire endpoint state,
spacing residue) determine the endpoint at every height $K\ge4$.
\emph{Proof.} Induct on the layer number using Lemma 3. The lower spacing
bound persists and the exact closed pair determines the next six
transfers. Merely seeing the visible state repeat at two sample powers
would not suffice.

The explicit entrance configurations and their cycles are listed in Appendix B. The conclusion is a conditional half-line theorem: it controls these configurations at every height. The next two sections establish that the signed process actually enters them.

\section{The signed boundary rule, including all larger powers}
\label{sec:5}

Let $M$ be a large power of 7 and $C_M=(M-1)/2$. Suppose that on each side the entire used set is described by a boundary record with holes and current coordinate within 512 of $C_M$. In these coordinates a cross-side jump of length $M$ is the reflection

\begin{equation*}
(s,x)\longmapsto(-s,1-x). \tag{4}\label{eq:reflection}
\end{equation*}

The finite boundary process offers all unused same-side power neighbors and the unused reflected point. It orders the candidates by their offset, preferring side + in a tie.

\noindent\textbf{Lemma 5 (signed boundary rule).} For $M\ge7^6$, with
the complete used-set records and the bounds just specified, the finite
boundary rule agrees with the signed greedy rule throughout each
certified boundary word.
\emph{Proof.} Any power below $M$ is at most $M/7$. It cannot cross zero from $C_M+x$ with $|x|<512$. If such a backward jump is too large to reach the bounded boundary region, it lands in the filled prefix; the corresponding forward jump is worse than the finite candidate supplied by the certificate. The jump $M$ is exactly (4). Any power at least $7M$ reaches a magnitude far larger than an available boundary candidate. The numerical margin $M>100000$ is more than sufficient for all these inequalities.

The separate checker scans offsets in increasing order, testing side + before side -, instead of using the generator's neighbor enumeration. It checks that every current point and chosen point, and each filled-front coordinate, stays within the stated bound.

\section{A uniform transition from scale $m$ to scale $7m$}
\label{sec:6}

\noindent\textbf{Proposition 6 (uniform scale transition).} Fix $k\ge6$
and the complete input state in Appendix A whose label is congruent to
$k$ modulo 36. Put $m=7^k$, $C=(m-1)/2$, and $D=(7m-1)/2=C+3m$.
Starting at index $m$ in that state, the exact greedy walk reaches the
next complete table state at index $7m$. The construction below gives
the transition and its count uniformly in $k$.

The certificate uses labels $40,\ldots,75$ for convenient representatives.
A label is a phase modulo 36, \emph{not} a lower bound on the exponent.
The proof is the following five-part construction; Section 7 verifies
phase closure and supplies the exact base from zero.

\subsection*{(i) Inner boundary}

Apply the boundary rule at $C$ until the front on the current side has offset greater than 100. These words have 99 to 111 choices. The other side's remaining used points have offsets below 100.

On the outgoing side, the half-line walk has a checked finite transient and then a translation cycle. For every transient and cycle current coordinate $x$, the checker explicitly verifies

\begin{equation*}
1-x<\text{the other side\textquotesingle s filled front}.
\end{equation*}

As the walk translates outward this remains true. Therefore a cross-side $m$ jump cannot reach an unused value. If subtracting $m$ no longer crosses zero, it instead lands in the filled prefix on the same side. Powers below $m$ are already covered by the exact half-line greedy checks. Until the first competing $7m$ jump, this outgoing walk is therefore exactly the half-line walk.

\subsection*{(ii) First crossing near $D$}

At cycle entry let the front be $f$ relative to $C$. Let $U_i$ and $W_i$ be the current and proposed next coordinates relative to that front in phase $i$, and let the cycle length be $P$. In cycle $q$ the opposite candidate from a $7m$ jump has offset

\begin{equation*}
6m+1-f-U_i-qP
\end{equation*}

relative to $C$. It is unused because the other side has no used points near $D$ yet. It defeats the same-side candidate $f+W_i+qP$ exactly when

\begin{equation*}
\begin{aligned}2qP&>6m+1-2f-U_i-W_i&&\text{if the current side is positive},\\2qP&\ge6m+1-2f-U_i-W_i&&\text{if the current side is negative}.\end{aligned}
\end{equation*}

Take the first pair ($q$,$i$) in chronological order satisfying this condition. Each phase threshold is an exact integer inequality. If $m$ is increased by $\Delta$ divisible by $P$, every threshold $q$ increases by $3\Delta/P$. Thus the winning phase and the entire resulting state relative to $D$ depend only on $m$ modulo $P$. No large numerical extrapolation is used here.

To justify the domain of this shift, the supplementary checker solves every phase threshold at the least $k\ge6$ in its exponent class. All thresholds are positive there (the smallest is 44105), and the winning phase and complete state agree with the certificate. Every threshold increases at larger scales in the class. The initial transient has current and chosen offsets below 2000, so its $7m$ reflection cannot win. The entire winning cycle has current and chosen magnitudes at most $D+9$; earlier cycles are farther inward. Hence the opposite crossing candidate is unused, and powers at least $49m$ are too large to compete.

The new point on the opposite side is $C+3m+A$ in magnitude, with $-3\le A\le3$, checked for all 36 phases. The first side is filled up to a bounded distance from $D$, and has only bounded extra used points there.

\subsection*{(iii) Return through the other side's old holes}

The next two choices in magnitude are $C+2m+A$ and $C+m+A$. They are unused and are better than any cross-side alternative. There are now marked endpoints at offsets $m+A$, $2m+A$, $3m+A$ on this side.

If $A$ was already used in the old inner boundary, descent from $m+A$ produces precisely the ladder of Section 4. The certified input is one of its 18 configurations; the finite bottom word and layer induction give the checkpoint immediately beyond $m+A$. The layer state at height $k$ is periodic in $k$ modulo 2 or 36.

If $A$ was unused, subtraction of $m$ from $m+A$ reaches $A$ directly. There is no high ladder. A checked finite word completes the old holes through $A$. Apply one marked-point transfer of spacing $m$ to reach a checkpoint past $m+A$. This exceptional case occurs only in the phase labeled 43 (equivalently $k\equiv7\pmod{36}$).

In either case apply a marked-point transfer of spacing $m$ to reach a checkpoint beyond $2m+A$. The point $3m+A$ remains marked.

These half-line operations are also valid signed operations. For the first two descending moves, the offered $m$ subtraction beats the possible $7m$ crossing by $m+1-2A$ and $3m+1-2A$, respectively. For the subsequent return use $H=4096$. Both current and chosen offsets from $C$ lie in $[-100,2m+H)$. The bottom computations and marker bounds give this for chosen points. The supplementary checker also bounds transient currents and verifies $\max(U)-W_i<2000$ for each marker cycle, as well as current bounds throughout each repair. Here $U$ denotes cycle current coordinates and $W_i$ the marked choice. Translation extends these bounds to every repeated cycle. The latest target is $2m+A$, and $A+2000<H$. The lower bound is checked through the finite bottom, transient and repair words; translation only increases the coordinates. The exceptional direct return has its own checked finite word. The accompanying reduction-premise checker records the corresponding current bounds.

A $7m$ crossing during this return has offset at least $4m+1-H$ and loses to the local candidate because $2m+1-2H>0$ for $m\ge7^6$. If an $m$ subtraction crosses zero, its opposite offset is $1-x\le101$, which is in the already filled first-side prefix since $3m-200>101$. Powers below $m$ cannot cross zero from this region's lower endpoint $C-100$, since their maximum is $m/7$. Higher powers are worse still: $49m>2(C+3m+H)$. Same-side backward moves into the filled prefix are excluded by the half-line rule. This verifies the signed/ half-line interface, rather than assuming the two processes agree.

\subsection*{(iv) Approach to the outer boundary}

Temporarily ignore the last marker $3m+A$. At the checkpoint beyond $2m+A$, take the checked transient to a translation cycle of length $P$. Write $f$ for the front relative to $2m+A$ after this transient. Execute

\begin{equation*}
q=\left\lfloor\frac{m-A-256-f}{P}\right\rfloor
\end{equation*}

complete cycles. The new front relative to $D$ is

\begin{equation*}
F=-256-\bigl((m-A-256-f)\bmod P\bigr). \tag{5}\label{eq:approach}
\end{equation*}

The supplementary checker verifies $q\ge0$ at the least exponent in each class (the smallest $q$ is 1586). Increasing $m$ by a multiple of $P$ increases $q$ by that multiple divided by $P$ and preserves $F$. The remainder identity alone would not establish that $q$ is a valid number of cycles.

The cycle certificate checks that its largest current or chosen coordinate, minus $P$, is less than 200. Thus every executed cycle's current and chosen offsets relative to $D$ are less than $F+200\le-56$. The transient is bounded by $2m+4096$ in offsets from $C$, also below $D-56$. Since $A\ge-3$, the already-used marker $D+A$ is not selected again. A $7m$ reflection of a current offset below -56 has offset above 57 and strictly loses to the offered local candidate. Thus ignoring the final marker was legitimate. Formula (5) and the remaining finite state depend only on $m$ modulo $P$. Restore the already-used point $D+A$ to the boundary record; this restoration does not count as a new visit.

\subsection*{(v) Outer boundary and the index}

Both sides now have entirely filled prefixes ending within 512 of $D$, with bounded additional used points. Use the boundary rule at scale $7m$. For records $(f_+,B_+),(f_-,B_-)$ relative to $D$, the number of used integers, including zero, is exactly

\begin{equation*}
7m-2+f_++f_-+|B_+|+|B_-|.
\end{equation*}

Therefore index $n=7m$ is reached exactly when

\begin{equation*}
f_++f_-+|B_+|+|B_-|=3. \tag{6}\label{eq:count}
\end{equation*}

The starting value in (6) is at most 3. Every greedy choice adds one to it. The separate checker follows exactly the required 256 to 276 choices, verifies all ties with positive preference, and obtains the next state in Appendix A. For the last phase the next state is the first phase. This counts actual visited integers; premarked points are never counted as additional steps.

\section{Closure for all exponents and coverage}
\label{sec:7}

All translation periods used in Section 6 divide $L=19240$. All layer cycles in Section 4 divide 36 and have stabilized by height 5. By (3), replacing $k$ by $k+36$ leaves $m$ modulo every required period unchanged. Equations (4)-(6), the first-crossing inequalities, and the endpoint transfer rule consequently give the \emph{same} normalized transition for every $k$ in a given phase. All size restrictions are preserved as $k$ increases. The common bound $m\ge7^6$ ensures the finitely many initial size restrictions.

The verifier checks 36 such transitions, with complete inputs and outputs, and checks that they close in order. This establishes a finite transition rule, not a comparison of two sampled orbits.

For the base case a separate uncompressed implementation generates every term from $n=0$ to $n=7^6=117649$. To bound its power search, let $R$ be the greatest magnitude already used and $p$ the first power exceeding $2R$. Then $\mathrm{current}+p$ is unused. For any larger power $q\ge7p$, $|\mathrm{current}\pm q|\ge q-R>p+R$. No such power can improve the available candidate. This proves the completeness of that finite neighbor enumeration. Its 117650 terms have the SHA-256 recorded in Section 9, and its endpoint is exactly phase 6 of the table. The supplied checker recomputes this prefix directly.

Induction now gives the table at every scale $k\ge6$. Its least filled-front offset is -36 on either side. Hence every magnitude at most $C_k-37$ is used, proving (1) and therefore coverage of $\mathbb Z$.

The initial run also checks phases 4 and 5 directly. Combining these finite checks with the induction yields period 36 for all $k\ge4$. No claim is needed here that 36 is the smallest possible period.

\section{Deviation bounds and the $3/8$ equality}
\label{sec:8}

\noindent\textbf{Corollary 7 (conditional deviation lower bound).} The uniform scale transition in Proposition 6 implies the lower bound $3/8$ in (9). In the Lean development, the formalized base and scale transitions now also imply the exact lower-bound epsilon statement directly from the original rule. The formal proof uses the phase-42 family described below.

The first-crossing argument below is conditional on the uniform transition
and does not constitute an independent proof of it. The all-index upper
bound and complete original-rule contract are established later in this
section by an additional formal assembly.

Write $E(n)=\bigl||a(n)|-n/2\bigr|$ and $m=7^k$, $k\ge6$. Let $n_k$ be the first crossing by a step of length $7m$ in the transition from scale $m$ to $7m$. The first-exit certificate gives

\begin{equation*}
n_k=4m+\delta_k,\qquad |a(n_k)|=\frac{7m-1}{2}+A_k. \tag{7}\label{eq:crossing}
\end{equation*}

where $-3\le A_k\le3$ and $-1\le\delta_k\le12$. Both offsets are functions of $k$ modulo 36. Here is the exact count establishing the index.

Immediately before the crossing, let $(F,B)$ describe the outgoing side relative to $D=(7m-1)/2$, with $B$ measured from its front, and let $(f_0,B_0)$ describe the other side relative to $C=(m-1)/2$, with $B_0$ measured from $C$. The two sides have

\begin{equation*}
D+F-1+|B|\qquad\text{and}\qquad C+f_0-1+|B_0|
\end{equation*}

used magnitudes. Include zero and the new crossing point, then subtract one to convert the number of used integers to the index. This gives

\begin{equation*}
\delta_k=F+f_0+|B|+|B_0|-2.
\end{equation*}

The new crossing point has offset $A_k$ from $D$. Substitution in $E$ gives the exact identity

\begin{equation*}
\frac{E(n_k)}{n_k}=\frac38+\frac{8A_k-4-7\delta_k}{8(4m+\delta_k)}. \tag{8}\label{eq:deviation}
\end{equation*}

The numerator of the correction is bounded, while $m$ tends to infinity. Hence $E(n_k)/n_k$ tends to $3/8$, and consequently

\begin{equation*}
\limsup_{n\to\infty}\frac{E(n)}n\ge\frac38. \tag{9}\label{eq:lower}
\end{equation*}

This is a proof of a subsequence limit from the scale-transition theorem, not an inference from a graph. It is not an upper bound for all intervening indices and does not establish equality in (9).

The supplementary \path{checker/check_crossing_corollary.py} derives all 36 offset pairs, checks the count by separately assembling the finite used sets, and checks (8) with rational arithmetic. A historical compressed-trace comparison recorded agreement for all 36 phases $k=4,\ldots,39$. This supplementary numerical comparison is not rerun by the supplied checker and is not needed for the corollary.

\subsection*{A machine-checked phase-42 proof}
Put $F(n)=8E(n)-3n$. For every integer $t\ge0$, the theorem
\path{original_rule_crossing42} gives the exact actual-walk identities
\begin{equation*}
k_t=6+36t,\quad m_t=7^{k_t},\quad n_t=4m_t+1,
\quad 2a(n_t)=-7m_t-3. \tag{10}\label{eq:phase42}
\end{equation*}
Here 42 is the certificate's phase label; its least admissible exponent
is 6. The actual checkpoint is obtained from the complete scale schedule.
The existing certified outgoing word then supplies both the elapsed count
and the signed coordinate in (10). Neither identity is an assumed
crossing formula or an extrapolation from finite simulations.

The theorem \path{original_rule_crossing_defect} proves the integer
identity in
\begin{equation*}
F(n_t)=5,\qquad
\frac{E(n_t)}{n_t}=\frac38+\frac{5}{8n_t}. \tag{11}\label{eq:phase42defect}
\end{equation*}
The ratio equality is its elementary rational rewriting. The indices are
unbounded: for any cutoff $N\ge0$, taking $t=N$ gives $n_t\ge N$.
Consequently \path{original_rule_frequent_lower} proves
\begin{equation*}
\forall p,q\in\mathbb N_{>0}\ \forall N\in\mathbb N\quad
\exists n\ge\max\{N,1\}:\quad -8pn\le qF(n). \tag{12}\label{eq:lowercontract}
\end{equation*}
For positive $n$ this is the rational-epsilon lower-limsup contract
$E(n)/n\ge3/8-p/q$ at arbitrarily late indices. Its only sequence
hypothesis is the literal \texttt{OriginalProblem.FollowsRule}; the
existing proofs of the base and transitions are imported into its proof.
The sources are \path{research/actual_lower_bound/}.

This formal route uses one sufficient infinite phase family. The general
36-phase subsequence-limit calculation in (7)--(9) remains a manuscript
derivation; the new Lean conclusion is the exact integer statement (12),
rather than a theorem expressed through a real-analysis library's
\texttt{limsup} interface.

\subsection*{Complete-state tail bounds}
Let $U_n=\{a(0),\ldots,a(n)\}$. For a nonnegative integer cutoff $D$, define
\begin{align*}
T_D(n)&=\#\{x\in U_n:xa(n)>0,\ |a(n)|<|x|\le D\},\\
Q_D(n)&=\#\{x\in U_n:|x|>D\}.
\end{align*}
Thus $T_D$ counts already visited points farther out on the current ray,
and $Q_D$ counts all exterior points on both rays. Exact counting gives
\begin{equation*}
n\le D+|a(n)|+T_D(n)+Q_D(n). \tag{13}\label{eq:physicalcount}
\end{equation*}
Indeed the central part is covered by at most $D+|a(n)|+1$ points;
the remaining points are in the tail or the exterior. The walk's $n+1$
visited points are distinct, including zero once. The Lean proof handles
the zero-current case as well. Moreover, the 36 complete checkpoint
records and actual-state reachability give $Q_{C_k}(n)\le7$ whenever
$k\ge6$ and $n\le7^k$. Inside the scale interval
$[7^k,7^{k+1})$, this exterior bound is applied at checkpoint $k+1$.
These results are in \path{research/physical_counting/}.

\noindent\textbf{Lemma 8 (complete signed-ladder tail bound).}
Choose one of the 18 configurations $(A,f,B)$ in Appendix B, put
$O=\{x:x<f\}\cup B$, and let $m=7^k$, $C=(m-1)/2$, $k\ge6$.
Suppose the actual walk at index $b$ has current point $s(C+A+m)$,
where $s\in\{-1,1\}$. In the ray coordinates of Section 2, its complete
used-offset set on side $s$ is exactly
\[
O\cup\{A+m,A+2m,A+3m\}.
\]
The opposite ray has a complete boundary record $(g,H)$ with
$g\ge3m-200$; zero is included once and ray coordinates represent only
positive magnitudes. Then the subsequent complete ladder word of
Proposition 4, including descent, bottom and ascent with the two outer
markers retained, has a finite length $\ell$ and satisfies
\begin{equation*}
T_D(b+i)\le55+6(k-4)
\quad(0\le i\le\ell, D\in\mathbb N). \tag{14}\label{eq:laddertail}
\end{equation*}
Its complete signed endpoint and exact move count agree with the checked
ladder construction.

\emph{Proof.} For a normalized local state $(c,B')$, the number of local
occupied offsets ahead of the current is bounded by
\[
G(c,B')=\max\{0,-c-1\}+\#\{x\in B':x>c\}.
\]
Kernel checks establish $G\le53$ at every required transfer state and
every prefix of all 18 complete bottom words, with equality of the
complete represented used sets. The full-transfer prefix theorem covers
the transient, arbitrarily many translated cycles, the partial cycle and
the repair, retaining every later marker. At a repair seam the target
is absorbed into the local state. Induction through the six transfers
per layer and through all $k-4$ layers therefore covers all remaining
markers by the original tower of $6(k-4)$ points. The forced descent fits
the same budget. Adding the two persistent outer markers gives (14).
The signed return envelope proves validity of the same constructed word
for both signs, and greedy uniqueness identifies all of its prefixes
with the actual walk from the stated complete input.

This argument is formalized by \texttt{actual\_ladder\_tail} in
\path{research/upper_ladder_prefix/ConcreteLadderPrefix.lean}.
The numerical local bound, tower cover and signed validity are proved
within that theorem. Its input-state match, scale radius and opposite-ray
filled bound remain explicit premises.

The subsequent \path{research/upper_phase_bridge/} development discharges
all these premises from the original rule for all 35 ordinary phase
families. It includes the two-move phase-43 core, whose tail is
at most four, in a single theorem for every $k\ge6$. This theorem retains
the exact actual start index, complete start and endpoint states, signed
validity and move count, and proves (14) at every prefix of the selected
return core. The division $k-6=36t+i$, $0\le i<36$, makes the phase
dispatch exhaustive, rather than a test of selected exponents.

The \path{research/upper_boundary_prefix/} development proves
$T_D\le39$ through every prefix of all 72 signed boundary words, given
their complete actual input states. For all 36 inner words it discharges
that input condition from the original rule at index $7^k$, for every
$k\ge6$. The subsequent return assembly also reaches the complete actual
outer inputs from the marked approach endpoints.

\subsection*{The all-index estimate and exact equality contract}
\noindent\textbf{Theorem 9 (original-rule deviation bound).}
Every sequence satisfying the rule of Section 1 obeys
\begin{equation*}
F(n)\le32771+6(k-4)
\quad(k\ge6,\ 7^k\le n\le7^{k+1}). \tag{15}\label{eq:uppergoal}
\end{equation*}
It satisfies the lower contract (12) and the upper contract
\begin{equation*}
\forall p,q\in\mathbb N_{>0}\quad\exists N\in\mathbb N\quad
\forall n\ge\max\{N,1\}:\quad qF(n)\le8pn. \tag{16}\label{eq:uppercontract}
\end{equation*}
These exact integer statements are the rational-epsilon formulation of
\begin{equation*}
\limsup_{n\to\infty}\frac{E(n)}n=\frac38. \tag{17}\label{eq:equality}
\end{equation*}

\emph{Proof of the scale estimate.} Fix $m=7^k$, $C=C_k$, $D=C_{k+1}$
and $L=k-4\ge2$. Before the first crossing, complete inner-boundary and
outgoing-cycle prefix descriptions give $F(n)\le16261$. The proof covers
every repetition and splits the offset index by quotient and remainder
of the exact cycle length. Its finite current, front and occupied-set
bounds refer to all 36 original schedule rows; there is no assumption
that an endpoint bound controls intervening states.

At the crossing and its two entrance moves, the old active ray lies
strictly below the current. The three marked entrance points therefore
give $T_D\le3$. Lemma 8 and its actual-input instantiation bound the
complete core by $55+6L$. The post-core transfers, including both
exceptional phase-43 transfers, have bound 56. The marked approach has
bound 44, and the outer boundary has bound 39. Full-state equalities
join these actual words. Their terminal state equals the next checkpoint;
injectivity identifies its index as exactly $7^{k+1}$. Splitting at each
word length, including both ends, proves
$T_D(n)\le55+6L$ at every index from the crossing through that checkpoint.

The radius bounds follow directly from the two complete checkpoints.
For $n>7^k$, the filled central interval and injectivity give
$|a(n)|>C-37$. For $n\le7^{k+1}$, inclusion in the next checkpoint's
entire used set gives $|a(n)|\le D+4096$; the finite offset bound is
checked on both rays of all 36 records. Hence throughout the return,
\[
C-4096\le|a(n)|\le D+4096,\qquad n\ge4m-1.
\]
The time bound follows from the exact crossing correction
$-1\le\delta_k\le12$ and the affine cycle-index relation at every exponent.
Combining these inequalities with (13) and $Q_D\le7$ gives
\[
8|a(n)|-7n\le32771,\qquad
n-8|a(n)|\le28737+6L.
\]
Their maximum is $F(n)$. This proves (15) on the return, while the
before-crossing bound covers the remaining indices. The phase division
$k-6=36t+i$ exhausts all $k\ge6$.

The formal theorem \path{original_rule_scale_bound} in
\path{research/upper_scale_assembly/Theorem.lean} proves (15) with only
the original rule and the displayed index bounds as hypotheses. The
exact discrete asymptotic argument in
\path{research/upper_asymptotic/FromScale.lean} then gives (16).
Together with (12), \path{original_rule_three_eighths} proves the named
target \texttt{OriginalRuleThreeEighths}. This is an inhabitant of that
proposition, not merely its definition or a conditional arithmetic lemma.
No additional transition, tail, reachability, time or radius premise
remains; the existing base and transition proofs are dependencies.
Equation (17) is the usual real-notation interpretation of the two
epsilon contracts; the Lean statement uses exact integers and does not
invoke a real-analysis library's \texttt{limsup} API.

\section{What was checked by a separate implementation}
\label{sec:9}

The generator enumerates power neighbors. The separate checker uses candidate-first scanning and repeated integer division for local and boundary moves; it imports no generator code. It checks:

\begin{itemize}
\item 119 uniform marked-endpoint transfer certificates;
\item all 18 bottom configurations and their closed layer maps;
\item all 36 signed scale transitions and their cyclic closure;
\item 2573 distinct exact local transitions;
\item 640839 literal local/boundary choices, including two direct marked-from-the-start simulations for every marker certificate;
\item a separate uncompressed initial run of 117650 signed terms;
\item the uniform locality inequalities and the residue identities.
\end{itemize}

The initial signed prefix has SHA-256 \nolinkurl{3a9478e37859b92ba9e2a5d3a3cd1576cbba5071ceed829baae5d32b777d0c2c} when encoded as little-endian signed 64-bit integers, including $a(0)$.

\begin{samepage}
Run from the root of the extracted repository:

\begin{verbatim}
python3 scripts/verify.py
\end{verbatim}
\end{samepage}

Python 3.9 or newer and its standard library suffice. Do not use Python's -O option: the repository checker explicitly refuses it because it disables assertion checks. Exploratory computations through $7^{75}$ provide supplementary numerical evidence; the base case used here requires only the 117650 terms specified above. The infinite conclusion depends on the uniform transition.

The command also checks file identity and negative tests. See \path{docs/verification.md} for the exact map from paper claims to checks. It additionally runs \path{checker/check_reduction_obligations.py}, which checks the finite premises added in the dated reduction audit: least-scale thresholds and cycle counts, anchor separation, current-coordinate bounds, and exact signed comparison margins. The original totals above are unchanged. The Python command checks finite premises; it does not itself prove the universal reduction. The separate Lean companion formalizes the exact base, complete scale transitions and coverage implication. In particular, the two representative tests in check\_first do not themselves establish its all-spacing claim; the algebraic phase-shift argument in Section 6(ii) is essential.

Regenerate the certificates without overwriting the committed copies:

\begin{verbatim}
python3 src/generate_certificate.py
\end{verbatim}

Reports are written under \path{build/}. The source package includes the fixed certificates and their checksums, together with the generation inputs and executed verification records. Reproducing the checks requires no earlier large-prefix simulation.

The 9 November research extension additionally checks local bounds for
all 119 transfers, eight signed-marker records and 36 approach records.
For the 18 complete bottom words, 442 shared bitmap chunks supply exact
full-state checks covering all 43172 prefixes, including their initial
states. These finite checks feed the universal prefix and ladder proofs
used in Lemma 8; they are not additional sampled large exponents.

\section{Code, certificates and reproducibility}
\label{sec:10}

The published source package \cite{package}, version 1.0.0, contains the
manuscript of 3 November 2026,
exact JSON certificates and generation inputs, the separate Python checker
and generator, Lean sources and toolchain pin, and verification evidence.
Its complete byte inventory is \path{release-manifest.json}. The file
\path{EXPORT_PROVENANCE.json} maps the distribution to its source inputs;
private development commit IDs in older records are historical provenance,
not revisions of the public repository.

The present manuscript belongs to the local release candidate 1.1.0,
dated 9 November 2026. Its added deviation proofs are not part of the
published 1.0.0 package. The historical version DOI
\begin{quote}\small
\mbox{\url{https://doi.org/10.5281/zenodo.23912682}}
\end{quote}
identifies 1.0.0 only; no DOI is assigned to this candidate. The project
repository is \url{https://github.com/gujord/heraclitus-base7}.
The candidate is self-contained and requires no private repository.

With Python and the pinned Lean 4.35.0 toolchain installed, run from the
extracted candidate root:
\begin{verbatim}
python3 scripts/check_release.py
python3 scripts/rebuild_release.py --jobs 4
\end{verbatim}
See \path{docs/reproducibility.md}. This runner checks the separate Python
certificates and regeneration, all 21 coverage source bindings, and freshly
compiles the entire 351-module project closure: 244 coverage modules and
107 research modules. It audits 4998 coverage declarations and 1595 research
theorem statements and their transitive axiom lists, and requires an actual
inhabitant of the exact equality target. It uses no prior compiled project
objects. Only the pinned compiler's standard library is prebuilt.

The controls include the 14 Python rejection/premise tests, 14 coverage
mathematical rejections and five receipt mutations, followed by five research
controls: a proof placeholder, false arithmetic, an incorrect crossing index,
a false actual crossing defect and an injected axiom. The mathematical false
claims also have their negations proved. The allowed transitive axioms are
\texttt{propext}, \texttt{Classical.choice} and \texttt{Quot.sound}.
The final candidate records its fresh execution under
\path{evidence/release-1.1.0/}; a preliminary export has no completed receipt.
Source and object identities are rechecked at completion.

Earlier development runs built the research extension in stages, reusing
identity-checked objects: first 75, then 94, then 107 research modules.
Those historical reports retain their original execution scopes. The
complete candidate rebuild is a same-project machine check, not independent
human review. The exact rational-epsilon equality contract is proved in Lean;
a real-analysis-library \texttt{limsup} conversion is not claimed.

For direct comparison with the mathematical definition, the following is
\texttt{FollowsRule} verbatim from \path{formal/Heraclitus/OriginalProblem.lean},
in namespace \path{Heraclitus.OriginalProblem}:
\begin{lstlisting}
def FollowsRule (a : Nat → Int) : Prop :=
  a 0=0 ∧ ∀ n : Nat,
    (¬∃ i : Nat, i≤n ∧ a i=a (n+1)) ∧
    (∃ j : Nat, (a (n+1)-a n).natAbs=7^j) ∧
    ∀ v : Int, (¬∃ i : Nat, i≤n ∧ a i=v) →
      (∃ j : Nat, (v-a n).natAbs=7^j) →
      (a (n+1)).natAbs≤v.natAbs ∧
      ((a (n+1)).natAbs=v.natAbs → v≤a (n+1))
\end{lstlisting}
The coverage theorem's declaration is:
\begin{lstlisting}
theorem coverage (a : Nat → Int) (rule : FollowsRule a) :
    ∀ x : Int, ∃ i : Nat, a i=x ∧ ∀ j : Nat, a j=x → j=i
\end{lstlisting}
Here \texttt{Nat} contains zero, so the allowed exponent includes $j=0$.
The last inequality gives the larger integer, hence the positive integer,
priority when absolute values tie. The theorem
\texttt{canonical\_follows\_rule} proves that the canonical sequence
satisfies this rule. Statement fidelity and the universal reduction in
Sections 3--7 remain the main targets for mathematical review.

Code is licensed under MIT; the manuscript, documentation and certificate
data are licensed under CC BY 4.0. The component boundaries and full texts
are given in \path{LICENSE} and \path{LICENSES/}.

\section{Generalization and limitations}
\label{sec:11}

The marked-point argument is not specific to seven: for any fixed base, a certified translation cycle gives a transfer depending on the marker spacing modulo its period. For an odd base $b$, the coordinates $C=(b^k-1)/2$ are integral, the boundary reflection is again $x\mapsto1-x$, the high descent creates $b-1$ endpoints per layer, and the two boundary scales differ by $(b-1)b^k/2$. Thus the proof architecture extends naturally to other odd bases.

What does \emph{not} automatically extend is the existence of a finite closed catalogue of entrance states and marker transfers. In an unrestricted experiment, inserting arbitrary extra markers into all discovered base-seven cycles generated many additional cycles and did not close before the imposed 200000-state limit. The proof uses the 119 transfers forced by its particular ladders and scale states, not arbitrary marker patterns. This restriction must be retained in any generalization.

An auxiliary check shows that a single premarked point in an otherwise clean base-seven half-line walk with incoming distance 1 through 7 always allows eventual local coverage. Its reduction has 22 marker phases and 690 normalized states, ending in cycles of lengths 1,4,8,74,130. This explains the earlier appearances of 74 and 130, but is not itself the global signed proof.

The priority for another base should be to identify its reachable marker
transfers and prove a finite closure, or replace that closure by a uniform
hole bound. Additional large prefixes alone do not establish it. The
all-index estimate (15) and the exact $3/8$ contract are proved here for
base seven; their numerical bounds and finite catalogues are not asserted
for other bases. Minimality of the 36-scale period remains unproved.

\section{Earlier work and possible novelty}
\label{sec:12}

The definition and the coverage problem belong to the existing Heraclitus-transform literature and OEIS work. OEIS A379719 gives coverage and a block description for powers of two \cite{ref1}; OEIS A383442 asks a broader coverage question \cite{ref2}. Sloane's 2025 slides discuss a general infinite-step-set/gcd criterion [4, slide 8]. In his Rutgers Experimental Math Seminar on 10 September 2026, he formulates the conjecture that the Heraclitus transform of a set $S$ of positive step sizes visits every integer if and only if $S$ is infinite and $\gcd(S)=1$ \cite{talk2026}. The 2026 slides also express doubt about sufficiency [8, slide 9]. Here $S=\{7^j:j\ge0\}$ is infinite and contains 1, so it satisfies both conditions. The present coverage argument addresses this particular set; it provides no conclusion for arbitrary step sets. None of these general ideas is claimed as new here.

OEIS A377091 announces, on 29 October 2026, an AI-assisted proof for square steps under review \cite{ref3}. That announcement does not describe a powers-of-seven theorem. No manuscript for that proof was obtained during this investigation, so possible overlap with Sloane, Xausa, Karlsson or collaborators cannot be ruled out.

Searches on 2 November 2026 covered the transform's name, powers of seven, odd bases, the periods 74 and 130, the signed initial prefix, the relevant OEIS entries, and the authors' names. Sloane's home page and publication list were also inspected \cite{ref5,ref6}, and the relevant 2025 slide was read and visually checked. No matching prior base-seven coverage result was found in the sources searched. Possible overlap with unpublished work remains unresolved. Search engines returned much unrelated philosophical material; this is a limitation, not positive evidence of novelty. No absence-of-results search proves priority.

A bounded follow-up search on 3 November 2026 queried the OEIS JSON API
for the signed initial prefix and for the phrases ``not yet in the
sequence'' together with ``power of 3'', ``power of 5'' or ``power of 7''.
These queries returned no matches. Exact queries, times, returned counts and search limitations are
recorded in the source package. The phrase search is sensitive to wording
and does not exclude results under different descriptions or unpublished
overlap.

The possible new content is therefore the specific base-seven coverage theorem, its exact 36-scale state law, and the explicit closed marked-endpoint mechanism. These are candidates for novelty, subject to mathematical review and comparison with unpublished work.

\begin{samepage}
\section*{Verification status and AI assistance}
The argument, certificate generator, separate checker and exposition were
developed using OpenAI Codex under Roger Gujord's direction. This is a
personal research project; use of Codex does not imply review or endorsement
by OpenAI. The generator and checker use different implementations, but
were developed within the same AI-assisted project.
\end{samepage}

The Python checks establish finite certificate claims. The Lean companion
supplies the complete exact base from zero, every successive scale
transition with full used-set endpoints and exact counts, and coverage for
any sequence satisfying the literal original rule. The definition includes
zero initially, unused destinations, every distance $7^j$ for $j\ge0$,
least absolute value and positive preference in ties. A canonical sequence
satisfying that definition is also proved to exist. Neither coverage nor
periodicity is assumed in the rule.

The recorded fresh build rebuilt 244 project modules, checked 21 source
bindings and audited 4998 declarations. It used Lean 4.35.0 with its standard
library. The recorded allowed axioms are \texttt{propext},
\texttt{Classical.choice} and \texttt{Quot.sound}. The trust boundary includes
the compiler and kernel and, separately, correspondence of the formal
statements to the intended mathematical problem. Exporters of finite data
are not verified parsers; the kernel checks propositions about the literal
data they emit. Python results and file hashes are not Lean axioms.

The research extension proves (10), the integer identity $F(n_t)=5$ in
(11), the exact lower contract (12), the all-index estimate (15) and the
upper contract (16) from the original rule. The combined equality target
is now a proved Lean proposition, using exact rational-epsilon contracts.
The broader 36-phase subsequence-limit derivation in (7)--(9) is not the
formal route used for the lower contract, and the real notation in (17)
is not a claim to have formalized a real-analysis library's limsup API.

These checks were executed within this project. Independent human
mathematical review of Sections 3--7 and statement fidelity remains
outstanding, as does independently performed reproduction. Those review
tasks are distinct from the completed machine checks and do not prevent
further formalization. Coverage for all odd bases, minimality of the
period 36 and novelty are not established.

\clearpage
\begin{thebibliography}{8}
\bibitem{ref1} OEIS A379719, powers-of-two Heraclitus transform, R. Sigrist.
\par\mbox{\small\url{https://oeis.org/A379719}}

\bibitem{ref2} OEIS A383442, triangular-number Heraclitus transform, N. J. A. Sloane.
\par\mbox{\small\url{https://oeis.org/A383442}}

\bibitem{ref3} OEIS A377091, square-step Heraclitus transform, R. Sigrist and subsequent comments, including Sloane and Xausa's announcement.
\par\mbox{\small\url{https://oeis.org/A377091}}

\bibitem{ref4} N. J. A. Sloane, In the Footsteps of Heraclitus, REU Seminar, Rutgers University, 24 June 2025, slides, especially slide 8.
\par\mbox{\small\url{https://neilsloane.com/doc/REU2025.pdf}}

\bibitem{ref5} N. J. A. Sloane, home page and recent talks/papers.
\par\mbox{\small\url{https://neilsloane.com/}}

\bibitem{ref6} N. J. A. Sloane, publication list.
\par\mbox{\small\url{https://neilsloane.com/doc/pub.html}}

\bibitem{package} Roger Gujord, The base-seven Heraclitus transform:
manuscript and reproducible proof package, version 1.0.0, 2026.
\par\mbox{\small\url{https://doi.org/10.5281/zenodo.23912682}}
\par The accompanying source
archive includes a complete \path{release-manifest.json}.
\par\mbox{\small\url{https://github.com/gujord/heraclitus-base7}}

\bibitem{talk2026} N. J. A. Sloane, Highlights from the First 400,000
Sequences in the OEIS, Experimental Math Seminar, Rutgers University,
10 September 2026. Video, approximately 09:30--10:15; slides, page 9.
\par\mbox{\small\url{https://www.youtube.com/watch?v=HAHUrfhU8Eg}}
\par\mbox{\small\url{https://neilsloane.com/doc/EMSep2026.pdf}}

\end{thebibliography}

\clearpage
\input{tables.tex}
\end{document}
